\documentclass[12pt, aspectratio=169]{beamer}

\usepackage{fontspec}
\usepackage{amsmath, amssymb}
\usepackage{pifont}
\usepackage[backend=biber, citestyle=authoryear, block=ragged, maxnames=3, maxbibnames=10, doi=true]{biblatex}
\usepackage[abbreviations]{foreign}
\usepackage[super]{nth}
\usepackage{minted}
\usepackage{appendixnumberbeamer}

\usepackage{unicode-math}
\defaultfontfeatures{Script=Latin,Kerning=On,Ligatures={Common,Rare},Contextuals = Alternate}
\setmainfont[BoldFont={Fira Sans SemiBold}]{Fira Sans}
\setsansfont[BoldFont={Fira Sans SemiBold}]{Fira Sans}
\setmonofont[BoldFont={Fira Code SemiBold}, StylisticSet={2,3,5}, Scale=MatchLowercase]{Fira Code}
\setmathfont[Scale=MatchUppercase]{TeX Gyre DejaVu Math}
\setmathfont[range=up]{Fira Sans}
\setmathfont[range=it]{Fira Sans Italic}
\setmathfont[range=bfup]{Fira Sans SemiBold}
\setmathfont[range=bfit]{Fira Sans SemiBold Italic}

%%%% Beamer Settings
\usecolortheme{seahorse}
\useinnertheme{rounded}
%\useoutertheme{shadow}

\setbeamerfont{frametitle}{size=\normalsize}
\setbeamerfont{framesubtitle}{size=\small}

\setbeamertemplate{blocks}[default]
\setbeamerfont{block title}{size=\normalsize}
\setbeamercolor{block title}{fg=black,bg=blue!30}
\setbeamercolor{block body}{fg=black,bg=blue!10}

\setbeamertemplate{footline}[frame number]
%\setbeamertemplate{footline}{}
\setbeamercolor{footnote}{fg=_note}
\setbeamercolor{footnote mark}{fg=_note}
\setbeamerfont{footnote}{size=\small}

\setbeamertemplate{navigation symbols}{}
\setbeamertemplate{itemize items}[triangle]
\setbeamertemplate{enumerate items}[default]
\setbeamersize{description width=1cm}
\setbeamersize{text margin left=5mm, text margin right=5mm}
\setbeamercovered{transparent}

%% \setbeamertemplate{note page}[plain]
%% \setbeamerfont{note page}{size=\footnotesize}
%% \addtobeamertemplate{note page}{\setbeamerfont{itemize/enumerate subbody}{size=\footnotesize}}{}
%% \setbeameroption{show notes on second screen=bottom}

% https://tex.stackexchange.com/questions/232168/normal-text-is-invisible-when-using-beamer-with-notes-and-xelatex
\makeatletter 
\def\beamer@framenotesbegin{% at beginning of slide
  \usebeamercolor[fg]{normal text}
  \gdef\beamer@noteitems{}% 
  \gdef\beamer@notes{}% 
}
\makeatother

\colorlet{_note}{red!70!black}

% https://tex.stackexchange.com/questions/305039/hrulefill-with-customized-rule
\newcommand*{\doublerule}{\hrule width \hsize height 1pt \kern 0.5mm \hrule width \hsize height 2pt}
\newcommand\doublerulefill{\leavevmode\leaders\vbox{\hrule height .8pt width .1pt\kern1pt\hrule height .8pt}\hfill\kern0pt}

%%%% TikZ
\usepackage{tikz}
\usetikzlibrary{calc,trees,shapes.callouts}

% https://tex.stackexchange.com/questions/193289/tikz-in-beamer-subscope-does-not-inherit-visible-on-from-parent-scope
\tikzset{
  % use option [visible on=<+->] to uncover parts of a tikzpicture
  invisible/.style={opacity=0},
  visible on/.style={alt=#1{}{invisible}},
  alt/.code args={<#1>#2#3}{%
    \alt<#1>{\pgfkeysalso{#2}}{\pgfkeysalso{#3}} % \pgfkeysalso doesn't change the path
  },
}

% Inline TikZ picture that behaves an anchored text for TikZ Overlay
\newcommand{\anchored}[2]{%
  \tikz[remember picture, baseline=(#1.base)] \node[inner sep=0pt, outer sep=0pt] (#1) {#2};%
}

%%%% BibTeX

\addbibresource[]{../bibliography.bib}
\renewcommand*{\bibfont}{\scriptsize}

\let\OrigCite\cite
\renewcommand*\cite[2][]{%
  {%
    \color{_note}%
    [\if\relax\detokenize{#1}\relax
      \OrigCite{#2}%
    \else
      \OrigCite[#1]{#2}%
    \fi]
  }%
}

%%%% Minted

\usemintedstyle{default}
\newminted[rocqcode]{elpi.py:RocqLexer}{fontsize=\small,linenos=true,escapeinside=\#\#,texcomments}
\newmintinline[rocq]{elpi.py:RocqLexer}{escapeinside=\#\#}
\newminted[elpicode]{elpi.py:ElpiLexer}{fontsize=\small,linenos=true,escapeinside=\#\#,texcomments}
\newmintinline[elpi]{elpi.py:ElpiLexer}{escapeinside=\#\#}

%%%% Math

\newtheorem{metatheorem}[theorem]{Metatheorem}

%%%% Terms

\usepackage{xspace}

\newcommand{\Coq}{{\scshape Coq}\xspace}
\newcommand{\Rocq}{{\scshape Rocq}\xspace}
\newcommand{\RocqElpi}{{\scshape Rocq-Elpi}\xspace}
\newcommand{\Elpi}{{\scshape Elpi}\xspace}

\newcommand{\cmark}{\ding{51}}%
\newcommand{\xmark}{\ding{55}}%

%%%% Title
\title{Tactics for equational reasoning in category theory}
\subtitle{Automating reasoning modulo associativity, functoriality, and cancellation}
\author{Kazuhiko Sakaguchi}
\institute{CNRS, ENS de Lyon, UCBL, LIP, UMR 5668, France}
\date{March 11, 2026}

\begin{document}

\setlength{\abovedisplayskip}{1ex}
\setlength{\belowdisplayskip}{1ex}

\begin{frame}[plain]
  \maketitle
\end{frame}

\begin{frame}[t]{Example (associativity)}

  \begin{minipage}[t]{.5\textwidth}
    \only<2->{\rocq{Fail rewrite H.}}

    \phantom{x}

    \only<3->{\rocq{rewrite compoA.}}

    \only<4->{\rocq{rewrite H.}}
  \end{minipage}%
  \begin{minipage}[t]{.5\textwidth}
    $H : g; h = u$

    \doublerulefill
    
    \only<-2>{$\visible<2>{\alert(}f; g\visible<2>{\alert)}; h = \dots$}%
    \only<3>{$f; \alert{(}g; h\alert{)} = \dots$}%
    \only<4>{$f; \alert{u} = \dots$}%
  \end{minipage}

\end{frame}

\begin{frame}[t]{Example (functoriality)}

  \begin{minipage}[t]{.5\textwidth}
    \only<1>{\phantom{x}}%
    \only<2->{\rocq{rewrite -Fcomp.}}

    \only<3->{\rocq{rewrite compoA.}}

    \only<4->{\rocq{rewrite Fcomp.}}

    \only<5->{\rocq{rewrite H.}}
  \end{minipage}%
  \begin{minipage}[t]{.5\textwidth}
    $F : \mathcal{C} \to \mathcal{D},$

    $H : F(g; h) = u$

    \doublerulefill
    
    \only<1>{$F(f; g); F(h) = \dots$}%
    \only<2>{$F((f; g); h) = \dots$}%
    \only<3>{$F(f; \alert{(}g; h\alert{)}) = \dots$}%
    \only<4>{$F(f); F(g; h) = \dots$}%
    \only<5>{$F(f); \alert{u} = \dots$}%
  \end{minipage}

\end{frame}

\begin{frame}[t]{Example (cancellation)}

  \begin{minipage}[t]{.5\textwidth}
    \only<1>{\phantom{x}}%
    \only<2->{\rocq{rewrite -[g]comp1o.}}

    \only<3->{\rocq{rewrite -(isoK f).}}

    \only<4->{\rocq{rewrite compoA.}}

    \only<5->{\rocq{rewrite H.}}
  \end{minipage}%
  \begin{minipage}[t]{.5\textwidth}
    $f : \mathrm{Iso}(a, b),$

    $g : \mathrm{Hom}(b, c),$

    $h : \mathrm{Hom}(a, c),$

    $H : f; g = h$

    \doublerulefill
    
    \only<1>{$g = f^{-1}; h$}%
    \only<2>{$\alert{1_b}; g = f^{-1}; h$}%
    \only<3>{$(\alert{f^{-1}; f}); g = f^{-1}; h$}%
    \only<4>{$f^{-1}; \alert{(}f; g\alert{)} = f^{-1}; h$}%
    \only<5>{$f^{-1}; \alert{h} = f^{-1}; h$}%
  \end{minipage}

\end{frame}

\begin{frame}{Solution: \rocq{reflexivity} and \rocq{rewrite} modulo the equational laws}

  \begin{itemize}
  \item \rocq{rewrite H} should suffice in any of the examples.
  \item The main differences with AAC Tactics~\cite{DBLP:conf/cpp/BraibantP11}:
    \begin{itemize}
    \item[\color{green}\cmark] dependently-typed morphisms,
    \item[\color{green}\cmark] cancellation (split monomorphisms/epimorphisms),
    \item[\color{red}\xmark] uninterpreted function symbols,
    \item[\color{red}\xmark] commutativity.
    \end{itemize}
    \medskip
  \item There is no support for functoriality yet.
    It requires a few more weeks of work, but should not be difficult.
  \end{itemize}

\end{frame}

\begin{frame}{Methodology: proof by reflection~\cite{DBLP:conf/lics/AllenCHA90}}

  \begin{tikzpicture}[semithick]
    \begin{scope}[every node/.style = {draw, rounded corners=1mm}]
      \node (input) at  (-5.5, 2) {input\vphantom{fg}};
      \node (tree)  at  (0,    2) {syntax tree\vphantom{fg}};
      \node (nf)    at  (5.5,  2) {normal form\vphantom{fg}};
    \end{scope}
    \draw[->, dashed] (input) to [bend right=10] node[below] {\vphantom{fg}reification} (tree);
    \draw[->]         (tree)  to [bend right=10] node[above] {\vphantom{fg}interpretation} (input);
    \draw[->]         (tree)  to [bend left=10]  node[above] {\vphantom{fg}normalisation} (nf);
    \node at (-5.5, 0) {$(a ; (\mathrm{id} ; b)) ; (c ; d)$};
    \node[inner sep=0, outer sep=0] at (0, 1) {}
      [child anchor=north,
       level/.style = {sibling distance={5cm/pow(2,####1)}, level distance=.4cm},
       every node/.style = {circle, fill=black!20, below, inner sep=-2pt, minimum width=1.2em}]
      child{
        child{node{$a$}}
        child{child{node{\textsf{id}}} child{node{$b$}}}}
      child{
        child{node{$c$}}
        child{node{$d$}}};
    \node at (5.5, 0) {$\lambda k.\, a ; (b ; (c ; (d ; k)))$};
    \begin{scope}[visible on=<2->]
      \node (yoneda) at (0, 4) {Yoneda translation~\cite{DBLP:conf/lics/JaberLPST16}};
      \draw[->,thick,red] (input) |- (yoneda.south) -| (nf);
    \end{scope}
  \end{tikzpicture}

\end{frame}

\begin{frame}{The Yoneda translation~\cite{DBLP:conf/lics/JaberLPST16}}

  \small

  \begin{definition}[Yoneda translation]
    Given a category $\mathcal{C} \coloneq (\mathrm{Ob}, \mathrm{Hom}, \mathrm{id}, ;)$,
    we define its Yoneda translation
    $\mathcal{C}_{\mathcal{Y}} \coloneq
     (\mathrm{Ob}_{\mathcal{Y}}, \mathrm{Hom}_{\mathcal{Y}},
      \mathrm{id}_{\mathcal{Y}}, ;_{\mathcal{Y}})$
    as follows.
    \begin{align*}
      \mathrm{Ob}_{\mathcal{Y}}
      &\coloneq \mathrm{Ob}, \\
      \mathrm{Hom}_{\mathcal{Y}}(p, q)
      &\coloneq \Pi r : \mathrm{Ob},\, \mathrm{Hom}(q, r) \to \mathrm{Hom}(p, r), \\
      \mathrm{id}_{\mathcal{Y}}(p)
      &\coloneq \lambda (q : \mathrm{Ob}) \, (k : \mathrm{Hom}(p, q)),\, k, \\
      f ;_{\mathcal{Y}} g
      &\coloneq \lambda (s : \mathrm{Ob}) \, (k : \mathrm{Hom}(r, s)),\, f \, s \, (g \, s \, k).
      & \left(\genfrac{}{}{0pt}{}{f : \mathrm{Hom}_{\mathcal{Y}}(p, q)}{g : \mathrm{Hom}_{\mathcal{Y}}(q, r)}\right)
    \end{align*}
  \end{definition}

  \begin{metatheorem}[Yoneda lemma]
    The Yoneda translation of a category is a category with laws that holds definitionally.
  \end{metatheorem}

\end{frame}

\begin{frame}{The Yoneda translation extended with definitional functoriality}

  \begin{itemize}
  \item From the PL perspective, the Yoneda translation is \emph{CPS translation}.
  \item Given a functor $F : \mathcal{C} \to \mathcal{D}$, we can embed $\mathcal{C}$ in the Yoneda
    translation of $\mathcal{D}$.
    \begin{itemize}
    \item $F$ can be seen as another piece of continuation.
    \end{itemize}
  \item Given there is a translation rule for functor applications,
    \begin{itemize}
    \item identity,
    \item associativity, and
    \item functoriality
    \end{itemize}
    should become definitional.
  \end{itemize}

\end{frame}

\begin{frame}{Implementation}

  \begin{itemize}
  \item Built for the Graph Rewriting library~\cite{DBLP:conf/cpp/ArsacHP26}
  \item 195 lines of code of \Rocq:
    \begin{itemize}
    \item Syntax tree,
    \item Interpretation and normalisation functions, and
    \item Soundness lemma,

      \ie, for any syntax tree, its normal form equals its interpretation.
    \end{itemize}
  \item 199 lines of code of \Elpi~\cite{Tassi:habilitation}:
    \begin{itemize}
    \item Reification, and
    \item Matching modulo AU/G.
    \end{itemize}
  \end{itemize}

\end{frame}

\begin{frame}{Miscellaneous}

  Reification and matching procedures make a good use of \Elpi:
  \begin{itemize}
  \item \emph{Constraints} made the interaction between reification and matching easier.
  \item The matching procedures use \emph{backtracking} extensively.
  \end{itemize}

  \bigskip

  Our tactics modulo cancellation are built around the claim that such reasoning amounts to
  inserting and deleting $I$:
  \begin{align*}
    I
    &\coloneq f; \, I^*; \, f^{-1}
    & \text{($f$ is a split monomorphism, \ie, $f; f^{-1} = \mathrm{id}$)} \\
    I^*
    &\coloneq \mathrm{id} \mid I; \, I^*.
    & \text{(The Kleene star on $I$)}
  \end{align*}
  But, the completeness of our method is still an \emph{open question}.

\end{frame}

\appendix

\begin{frame}[allowframebreaks]{References}
 \printbibliography
\end{frame}

\end{document}
